\section{Results}
\label{sec:results}

% NOT WRITTEN. The production run has not been executed. Nothing here may be
% drafted from pilot output: the pilot has been looked at, which is what a
% preregistered protocol exists to prevent for reported results.
%
% Every number must be \input from a generated file. If a value cannot be
% traced through replication, seed, scenario, estimator, metric, gen_surv
% version and commit, it does not belong in the paper.
%
% Planned structure:
%
% \subsection{Does discrimination track recovery of the truth?}
%   Figure 1. The headline is the vertical spread at a fixed C-index, not the
%   correlation coefficient: report the range of MISE consistent with an
%   apparently acceptable concordance.
%
% \subsection{Which conventional metrics detect the error?}
%   Figure 2 and the calibration measures. This is where the contribution
%   lands: for each conventional metric, how much truth-error can hide behind
%   an acceptable value of it.
%
% \subsection{Where misspecification bites}
%   Figures 3-5: MISE against effect size, censoring and sample size, with
%   Monte Carlo error bars throughout.
%
% \subsection{Parameter recovery is not distribution recovery}
%   Table 4, for the three (mechanism, estimator) pairs where a fitted
%   coefficient estimates a generating parameter. The point is that beta-hat
%   can be unbiased while S-hat is not close.
%
% \subsection{The adequacy region}
%   Figure 6, across a range of epsilon. Do not present one epsilon as a
%   threshold.
%
% \subsection{Failures}
%   Table 5 and Figure 8: P(fit failure | scenario, estimator).
